\documentclass[10pt]{article}

% Ћирилица
\usepackage[T2A]{fontenc}
\usepackage[utf8]{inputenc}
\usepackage[serbianc]{babel}
\usepackage{tabularx}

\usepackage[a4paper,top=1.5cm,bottom=1.5cm,left=1.8cm,right=1.8cm]{geometry}
\usepackage{amsmath}
\usepackage{multicol}
\usepackage{titlesec}
\titleformat*{\section}{\large\bfseries}
\titlespacing*{\section}{0pt}{6pt}{3pt}
\usepackage{enumitem}
\usepackage{tcolorbox}
\tcbuselibrary{breakable}
\usepackage[european]{circuitikz}
\usetikzlibrary{arrows.meta, babel}

\usepackage{diagbox}
\usepackage{colortbl} % (опционо, ако желиш сиву позадину заглавља као у књизи)

\pagestyle{empty}
\setlength{\columnsep}{0.8cm}
\setlength{\parindent}{0pt}
\setlength{\parskip}{3pt}

% Оквир за кључне податке
\newtcolorbox{kljucno}{colback=blue!4, colframe=blue!45!black,
    boxrule=0.5pt, arc=1.5pt, left=6pt, right=6pt, top=3pt, bottom=3pt}

% Решени пример
\newtcolorbox[auto counter]{primer}{
    colback=white, colframe=black!55, colbacktitle=black!55, coltitle=white,
    fonttitle=\bfseries\small, title={Пример~\thetcbcounter},
    boxrule=0.5pt, arc=1.5pt, left=5pt, right=5pt, top=2pt, bottom=2pt,
    before skip=5pt, after skip=5pt}

% Задатак за вежбу (решење се даје само као резултат)
\newtcolorbox[auto counter]{zadatak}{
    colback=green!3, colframe=green!40!black, colbacktitle=green!40!black, coltitle=white,
    fonttitle=\bfseries\small, title={Задатак~\thetcbcounter},
    boxrule=0.5pt, arc=1.5pt, left=5pt, right=5pt, top=2pt, bottom=2pt,
    before skip=5pt, after skip=5pt}

\begin{document}

\begin{center}
    {\LARGE\bfseries Лабораторијске вежбе из основа електротехнике} \\
    {\Large Мерење физичких појава.} \\
\end{center}
\vspace{10pt}

\noindent
\textbf{Име и презиме:} \hrulefill \quad
\textbf{Одељење:} \underline{\hspace{2cm}} \quad
\textbf{Група:} \underline{\hspace{1.5cm}}

\vspace{15pt}

\section{Вежба 1}

У овој вежби потребно је израчунати висине својих другара из одељења.
Висина се мери са адекватним инструментом, метром за мерење. Измерену висину
потребно је уписати у табелу, а онда мерење поновити још 3 пута. Потребно је тако
измерити 4 својих другара. На крају, потребно је израчунати средњу вредности $\tilde{x}$ и 
мерну несигурност мерења $u$.

\textbf{Кораци рада:}
\begin{enumerate}
  \item Одабрати 4 испитаника и уписати њихова имена 
  у прву колону табеле (\textbf{ID1}--\textbf{ID5}).
  \item Поставити испитаника да стоји усправно и помоћу метра измерити 
  његову висину у метрима ($\text{m}$).
  \item Измерену вредност уписати у колону за прво мерење.
  \item Поновити поступак мерења још 5 пута за истог испитаника (сваки пут из почетка мерити
  ) и резултате уписати у остале колоне.
  \item Поновити кораке 2-3-4 за преостала четири испитаника.
  \item На основу прикупљених података у табели, за сваког испитаника израчунати средњу 
  вредност висине $\tilde{x}$ и мерну несигурност резултата $u$ и израчунати финални резултат $x$.
  Упоредити одступања између појединачних мерења и правих резултата.
\end{enumerate}

% \begin{table*}[htbp]
%   \centering
%   \renewcommand{\arraystretch}{1.8}
%   \setlength{\tabcolsep}{4pt} 
%   \begin{tabular}{|l|c|c|c|c|c|c|c|c|}
%     \hline
%     \textbf{Име особе} \hspace{0.8 cm} & \textbf{1} & \textbf{2} & \textbf{3} & \textbf{4} &  $\tilde{x}$ & $u$ & \textbf{Резултат $x$}  \\ \hline
%     & \hspace{1.1cm} & \hspace{1.1cm} & \hspace{1.1cm} & \hspace{1.1cm} & \hspace{0.8 cm} &  \hspace{0.8 cm} &  \hspace{1.6cm} \\ \hline
%     & & & & &  & & \\ \hline
%     & & & & & & & \\ \hline
%     & & & & & & & \\ \hline
%     & & & & & & & \\ \hline
%   \end{tabular}
% \end{table*}

\begin{table*}[htbp]
  \centering
  \renewcommand{\arraystretch}{1.8}
  \begin{tabularx}{\textwidth}{|p{2.8cm}|>{\centering\arraybackslash}X|>{\centering\arraybackslash}X|>{\centering\arraybackslash}X|>{\centering\arraybackslash}X|>{\centering\arraybackslash}X|>{\centering\arraybackslash}X|c|}
    \hline
    \hline
    \textbf{Име особе} \hspace{0.8 cm} & \textbf{1} & \textbf{2} & \textbf{3} & \textbf{4} &  $\tilde{x}$ & $u$ & \textbf{Резултат $x$}  \\ \hline
    & \hspace{1.1cm} & \hspace{1.1cm} & \hspace{1.1cm} & \hspace{1.1cm} & \hspace{0.8 cm} &  \hspace{0.8 cm} &  \hspace{1.6cm} \\ \hline
    & & & & &  & & \\ \hline
    & & & & & & & \\ \hline
    & & & & & & & \\ \hline
    & & & & & & & \\ \hline
  \end{tabularx}
\end{table*}

\newpage
\section{Вежба 2}

У овој вежbи потребно је измерити електричне отпорности себе и предмета око себе. 
Пре свега потребно је измерити отпорност себе, поновити мерење 3 пута, и уписати у табелу. 
Изабрати, из своје околине, 3 предмета и поновити иста мерења и за њих. На крају, потребно 
је израчунати средњу вредности $\tilde{x}$ и 
мерну несигурност мерења $u$.


\textbf{Кораци рада:}
\begin{enumerate}
  \item Подесити мултиметар за мерење електричне отпорности ($\Omega$). Поставити преклопник на 
  највећи мерни опсег ($1 M \Omega$).
  \item Ухватити врхове обе игле за мерење прстима леве и десне руке (од прста до прста) и 
  очитати вредност отпорности са казаљке или екрана. Измерену вредност у омима $\Omega$ уписати у табелу.
  \item Измерити отпорност 4 различита предмета у лабораторији 
  (нпр. графитна мина оловке, метални кључ, пластични лењир, дрвени сто, 
  гумица за брисање). Уколико се казаљка не помера пуно, или екран показује баш мали број(много нула),
  то значи да је изабран погрешан мерни опсег. Потребно је поставити преклопник на мањи мерни опсег отпора. 
  \item На основу прикупљених података у табели, за сваки предмет или особу израчунати средњу 
  вредност отпорности $\tilde{x}$ и мерну несигурност резултата $u$ и израчунати финални резултат $x$.
  Упоредити одступања између појединачних мерења и правих резултата. Приметити различитост у отпорностима. Да ли је предмет изоратор или диелектрик?
\end{enumerate}

% \begin{table*}[htbp]
%   \centering
%   \renewcommand{\arraystretch}{1.8}
%   \setlength{\tabcolsep}{4pt} 
%   \begin{tabular}{|l|c|c|c|c|c|c|c|c|}
%     \hline
%     \textbf{Назив} \hspace{0.8 cm} & \textbf{1} & \textbf{2} & \textbf{3} & \textbf{4} &  $\tilde{x}$ & $u$ & \textbf{Резултат $x$}  \\ \hline
%     & \hspace{1.1cm} & \hspace{1.1cm} & \hspace{1.1cm} & \hspace{1.1cm} & \hspace{0.8 cm} &  \hspace{0.8 cm} &  \hspace{1.6cm} \\ \hline
%     & & & & &  & & \\ \hline
%     & & & & & & & \\ \hline
%     & & & & & & & \\ \hline
%     & & & & & & & \\ \hline
%   \end{tabular}
% \end{table*}
% У преамбулу додај:
% \usepackage{tabularx}

\begin{table*}[htbp]
  \centering
  \renewcommand{\arraystretch}{1.8}
  \begin{tabularx}{\textwidth}{|p{2.8cm}|>{\centering\arraybackslash}X|>{\centering\arraybackslash}X|>{\centering\arraybackslash}X|>{\centering\arraybackslash}X|>{\centering\arraybackslash}X|>{\centering\arraybackslash}X|c|}
    \hline
    \hline
    \textbf{Назив} & \textbf{1} & \textbf{2} & \textbf{3} & \textbf{4} & $\tilde{x}$ & $u$ & \textbf{Резултат $x$} \\ \hline
    & & & & & & & \hspace{2.5cm} \\ \hline
    & & & & & & & \\ \hline
    & & & & & & & \\ \hline
    & & & & & & & \\ \hline
    & & & & & & & \\ \hline
  \end{tabularx}
\end{table*}

\begin{multicols}{2}

\end{multicols}


\end{document}